\scanpage{51}
\mspage{49}
% lines 1-4
cas de $X$ géom. normal. \deleted{\unsure{Le résultat est} \unsure{$U\subset X$}} \inserted{On peut supposer de} plus $X$ géom. intègre (\deleted{\unsure{en raisonnant}} les composantes connexes l'étant). Si $U$ est un ouvert non vide de $X$, on a $H^1(X_{\overline K},\mathbb{Q}_\ell)\hookrightarrow H^1(U_{\overline K},\mathbb{Q}_\ell)$, donc on \deleted{peut}

% lines 5-13: large crossed block
\deleted{supposer $X$ affine, donc quasi-projectif. Soit $\hat X$ \unsure{est associé à} $U$. \unsure{Pour} $X$, \unsure{Et on peut donc} supposer $X$ affine, donc quasi-projectif. Quitte à faire encore une extension \inserted{finie} radicielle sur $K$, on peut supposer que $X\subset\overline X$, où $\overline X$ est projectif, géom. normal. Notons que \unsure{pour remplacer} $X$ par $U$. Prenant $U$ assez petit, $U$ est un revêtement étale fini d'un ouvert $V$ de $\mathbb{P}^n$. Alors on a $H^1(\unsure{\cdot}$}

% lines 14-20: further cancelled attempt, with additions
\deleted{\unsure{remplacer}} peut \deleted{\unsure{supposer} $X$ affine par $U$, qu'on peut supposer} \inserted{$X$} affine \inserted{et lisse}.

\deleted{Soit \unsure{Considérons} une clôture projective de \inserted{$X$} $U$, qu'on peut supposer (quitte à faire encore une extension radicielle $K'/K$, finie) géom. normale. Alors $H^1(U,\mathbb{Z}/\ell^\nu\mathbb{Z})$ classifie les \unsure{rev. galoisiens} \illegible{2} de $X$, \unsure{mais} \unsure{normaux} sur $X$, qui sont ramifiés seulement aux \unsure{pts} de $X-U$.}

% lines 21-29
Un procédé de sections transversales prises ds $V$ et de ``courbe générique'' nous ramène au cas où $X$ est de dim. 1. \deleted{Alors le calcul explicite \unsure{revient} \unsure{th.}} \inserted{De plus} \deleted{\unsure{soit}} supposons que le modèle projectif régulier $\hat X$ de $X$ est lisse. La dualité nous ramène au cas de $\unsure{H^1_!|H^1_c}(X_{\overline K},\mathbb{Q}_\ell)$, la suite exacte au cas de $H^1_!(\hat X,\mathbb{Q}_\ell)=H^1(\hat X,\mathbb{Q}_\ell)$. On est donc ramené au cas d'une \uline{courbe} \inserted{\uline{régulière}} \uline{lisse} sur $K$. À expliciter en

\uline{Proposition.} 9.8. \unsure{. . . . . .}
% archive sheet number: 50
